\section{A replacement word whose position ratios do not converge}
\label{sec:a289128}

Clark Kimberling's entries A289128--A289131, introduced on June 28,
2017, concern the left-to-right nonoverlapping replacement
\[
00\longmapsto0010,\qquad01\longmapsto011,\qquad10\longmapsto010.
\]
Let $R$ denote one pass of this scanner. A matched pair is consumed,
and newly emitted letters are not scanned during that pass. An
unmatched input letter is copied; in particular, an initial $11$
causes the first $1$ to be copied before scanning continues. Set
$w_0=00$ and $w_{n+1}=R(w_n)$. The first words are
\[
00,\quad0010,\quad0010010,\quad00100100110,
\]
followed by $001001001100101010$ and
$0010010011001010100110110110$.

The entry A289131 conjectures that the lengths satisfy its printed
order-five recurrence. The adjacent entries A289129 and A289130
conjecture that the positions of the $k$-th zero and one, divided by
$k$, tend respectively to $5/2$ and $5/3$. We prove the length
assertion and prove that neither position ratio has any real limit.
The repeated length assertion in A289128 is counted once.

\paragraph{Theorem.}
The literal words $w_n$ are nested prefixes and determine a unique
infinite binary word $W$. Write $L_n=|w_n|$ and let $Z_n,O_n$ be the
counts of its zeros and ones. For every $m\ge0$,
\[
\begin{array}{c|cc}
 & n=2m+1 & n=2m+2\\ \hline
L_n & 10\cdot2^m-3m-6 & 14\cdot2^m-3m-7\\
Z_n & 4\cdot2^m-1 & 6\cdot2^m-1\\
O_n & 6\cdot2^m-3m-5 & 8\cdot2^m-3m-6
\end{array}
\tag{\ref{sec:a289128}.1}
\]
The exceptional initial length is $L_0=2$. With the first six values
$2,4,7,11,18,28$, these lengths satisfy exactly
\[
L_n=L_{n-1}+3L_{n-2}-3L_{n-3}-2L_{n-4}+2L_{n-5}
\qquad(n\ge6).
\tag{\ref{sec:a289128}.2}
\]
Let $a(k)$ and $b(k)$ be the one-based positions of the $k$-th zero
and one in $W$. Both are defined for every $k\ge1$. Their ratios have
the subsequential limits
\[
\begin{aligned}
\frac{a(Z_{2m+1})}{Z_{2m+1}}&\longrightarrow\frac52,
&\frac{a(Z_{2m+2})}{Z_{2m+2}}&\longrightarrow\frac73,\\
\frac{b(O_{2m+1})}{O_{2m+1}}&\longrightarrow\frac53,
&\frac{b(O_{2m+2})}{O_{2m+2}}&\longrightarrow\frac74.
\end{aligned}
\tag{\ref{sec:a289128}.3}
\]
All four occurrence-rank subsequences are strictly increasing.
Consequently neither $a(k)/k$ nor $b(k)/k$ converges.

\paragraph{Proof: exact blocks and scanner boundaries.}
For $t\ge0$, define $A_t=0(10)^t$ and $B_t=0(110)^t$. For any finite
binary suffix $u$, the scanner satisfies
\[
R(A_tB_su)=B_tA_{2s+1}R(u),\qquad R(A_t)=B_t.
\tag{\ref{sec:a289128}.4}
\]
To verify the first equality, the initial $01$ pairs of $A_t$ emit
$(011)^t$ and leave one $0$. The first $0$ of $B_s$ combines with
that letter to form $00$, emitting $0010$. Each of the remaining
$s$ blocks $110$ is parsed as an unmatched $1$ followed by $10$,
and emits $1010$. Hence the output before $R(u)$ is
\[
(011)^t0010(10)^{2s}=0(110)^t0(10)^{2s+1}=B_tA_{2s+1}.
\]
The second equality follows from
$(011)^t0=0(110)^t$. These identities hold for arbitrary suffixes,
so they justify concatenating the particular blocks used below;
no concatenation property for arbitrary binary strings is assumed.

Put
\[
t_m=2^{m+1}-1,\qquad P_0=0010,\qquad
P_{m+1}=P_mA_{t_m}B_{t_m}.
\]
We prove simultaneously that
\[
R(P_mu)=R(P_m)R(u)\quad\hbox{for every }u,
\qquad R(P_m)=P_mA_{t_m}.
\tag{\ref{sec:a289128}.5}
\]
For $m=0$, the word $0010$ is consumed as the pairs $00,10$, and
$R(0010)=0010010=P_0A_1$. The scanner is at a complete boundary
after those input pairs. If (\ref{sec:a289128}.5) holds for $m$,
then (\ref{sec:a289128}.4) gives
\[
R(P_mA_{t_m}B_{t_m})
=P_mA_{t_m}B_{t_m}A_{2t_m+1}
=P_{m+1}A_{t_{m+1}},
\]
and the same computation with a suffix gives the boundary identity
at stage $m+1$. Therefore induction on the original iterates yields
\[
w_{2m+1}=P_m,\qquad w_{2m+2}=P_mA_{t_m}.
\tag{\ref{sec:a289128}.6}
\]
The next odd iterate appends $B_{t_m}$. Every word is thus a prefix
of its successor, including $w_0=00$.

\paragraph{Counting and the length recurrence.}
The block $A_t$ has length $2t+1$, with $t+1$ zeros and $t$ ones.
The block $B_t$ has length $3t+1$, with $t+1$ zeros and $2t$ ones.
Starting from $P_0=0010$, induction in
(\ref{sec:a289128}.6) proves all six formulas in
(\ref{sec:a289128}.1). In particular the lengths are unbounded,
so the nested prefixes determine $W$.

The length formulas for a fixed parity are a multiple of $2^m$
plus an affine function of $m$. They give the shorter relation
\[
L_{r+4}+2L_r=3L_{r+2}+3\qquad(r\ge1).
\]
Subtracting the relation at $r$ from the relation at $r+1$ gives
(\ref{sec:a289128}.2). Directly scanning the first five passes
gives its six initial values. This proves the conjectured equality
with the separately defined A289131 recurrence for every index.
Colin Barker already recorded the scalar parity formulas and
generating function in A289131 on July 2, 2017; the present argument
proves their equality with the original replacement lengths.

\paragraph{Infinite occurrence positions.}
Every positive-index iterate ends in $10$. Its final zero is
therefore the $Z_n$-th zero of the infinite word, and its final one
is the $O_n$-th one. Thus
\[
a(Z_n)=L_n,\qquad b(O_n)=L_n-1\qquad(n\ge1).
\tag{\ref{sec:a289128}.7}
\]
The count formulas show that the odd and even rank subsequences
increase strictly and tend to infinity. They also prove that both
symbols occur infinitely often. Substituting
(\ref{sec:a289128}.1) into (\ref{sec:a289128}.7) gives
\[
\begin{aligned}
\frac{a(Z_{2m+1})}{Z_{2m+1}}
 &=\frac{10\cdot2^m-3m-6}{4\cdot2^m-1},&
\frac{a(Z_{2m+2})}{Z_{2m+2}}
 &=\frac{14\cdot2^m-3m-7}{6\cdot2^m-1},\\
\frac{b(O_{2m+1})}{O_{2m+1}}
 &=\frac{10\cdot2^m-3m-7}{6\cdot2^m-3m-5},&
\frac{b(O_{2m+2})}{O_{2m+2}}
 &=\frac{14\cdot2^m-3m-8}{8\cdot2^m-3m-6}.
\end{aligned}
\]
Divide by $2^m$ and use $m/2^m\to0$ to obtain
(\ref{sec:a289128}.3). Were either full ratio convergent, both of
its displayed subsequences would have the same limit. Since
$5/2\ne7/3$ and $5/3\ne7/4$, neither converges. This refutes the
two proposed limits without relying on a finite-prefix
approximation. $\square$

\paragraph{Formal scope.}
Lean 4.19.0, with mathlib pinned at
\texttt{c44e0c8ee63ca166450922a373c7409c5d26b00b}, checks the entire
chain. The scanner uses actual lists of bits and recursively
consumes the original input. The exact block theorem is
\texttt{Kimberling289128.word\_blocks}. Theorems
\texttt{word\_nested}, \texttt{bit\_stabilizes}, and
\texttt{infinite\_word\_unique} establish the nested limit. The
infinite letter at zero-based position $i$ is constructed as the
$i$-th letter of $P_i$; the proof gives $|P_i|\ge i+4$ and
stabilization in every iterate of index at least $2i+1$.

The one-based position maps use \texttt{Nat.nth}. The theorem
\texttt{position\_spec} verifies their letters and exact preceding
occurrence counts at every positive rank. Theorems
\texttt{odd\_positions} and \texttt{even\_positions} prove
(\ref{sec:a289128}.7) for those maps. Four real-limit theorems and
strict monotonicity of the rank sequences imply
\texttt{position\_ratios\_do\_not\_converge}, which excludes every
real limit for both ratios. The source-specific endpoints are
\texttt{A289129\_limit\_false} and
\texttt{A289130\_limit\_false}. The length endpoint
\texttt{original\_length\_conjecture} identifies every literal
length with the independent integer recurrence, including the
exceptional zeroth word and exact recurrence start at index six.
All endpoint axiom dependencies are among
\texttt{propext}, \texttt{Classical.choice}, and \texttt{Quot.sound}.

\paragraph{Sources and prior-work screen.}
The package preserves the complete official OEIS exports for
A289128 revision 7, A289129 revision 4, A289130 revision 6, and
A289131 revision 19. Their saved Git blob hashes match the
official repository responses. The source review records searches
with both available engines, exact and semantic checks of the
current trureturing catalogue, and the existing mathematics
repository. No earlier proof of this exact length assertion or
refutation of these two exact position limits was located. This
bounded review does not certify historical priority.

The trureturing proof of A288476 uses the different rule
$00\to0101,1\to011$ and explicitly excludes the present position
and limiting-word questions. Michel Dekking's 2022 proof in
A288936 sends $10$ to $011$, rather than the $010$ used here, and
has different growth. The repository's earlier pair-rule and
morphic-discrepancy results also concern different substitutions.
Standard scanner recoding and subsequence methods are credited as
methods, not claimed as new general theory. The four limits above
are not asserted to classify all cluster points or determine exact
liminf and limsup values; corrected all-index inequalities are
outside this package's claimed scope.
